% TOP_PROBLEM: {"id": 2400013, "title": "Flint Hills series", "category_id": 8, "category": {"id": 8, "name": "analysis", "display_name": "Analysis", "description": "Limits, continuity, calculus, and function theory.", "slug": "analysis", "order_index": 8, "created_at": "2026-07-31T15:26:25.671Z"}, "status": "open"}
% FIRST_POSTED: 2026-10-06
% ATTEMPT_STATUS: solved
% OPENAI_CLAIM: {"schema_version":1,"source_repo":"https://github.com/openai/math","source_commit":"adc7f1241b42e322a6451854ab7e4b4c146bf78a","release_date":"2026-10-06","match":"full","resolution":"proved","independently_reviewed":false,"second_pass":"agreed","families":[{"family":"017","title":"The irrationality exponent of <i>π</i> is 2","subject":"Number theory","manuscripts":[{"title":"The irrationality exponent of π is 2","pdf_path":"preprints/The-irrationality-exponent-of-pi-is-2-September-24-2026/paper.pdf","readme_path":"preprints/The-irrationality-exponent-of-pi-is-2-September-24-2026/README.md","date":"2026-09-24","theorem_ref":"thm:flinthills (main.tex:90); thm:main (main.tex:64)","statement_sources":[{"path":"preprints/The-irrationality-exponent-of-pi-is-2-September-24-2026/build/main.tex","line":90,"label":"thm:flinthills"},{"path":"preprints/The-irrationality-exponent-of-pi-is-2-September-24-2026/build/main.tex","line":64,"label":"thm:main"}]}],"lean":{"doc_path":"lean/docs/017.md","comparators":["lean/ComparatorChallenges/PiExponent.json","lean/ComparatorChallenges/PiExponent.lean"],"declarations":[],"covers_main_theorem":true},"reasoning_trace":"reasoning_traces/irrationality-exponent-of-pi.pdf"}]}
\documentclass[11pt]{article}
\usepackage[margin=25mm]{geometry}
\usepackage{fontspec}
\setmainfont{Latin Modern Roman}
\usepackage{amsmath,amssymb,mathtools}
\usepackage{unicode-math}
\setmathfont{Latin Modern Math}
\usepackage{enumitem}
\usepackage{xurl,hyperref,bookmark}
\hypersetup{colorlinks=true,urlcolor=blue}
\setcounter{secnumdepth}{0}
\setlength{\parindent}{0pt}
\setlength{\parskip}{5pt}
\setlength{\emergencystretch}{3em}
\begin{document}
\section*{Flint Hills series}

\subsection{OpenAI's claim}
OpenAI presents the following claims relevant to this question, with manuscript numbers following paper order in Sources. In manuscript 1 (family 017; 2026-09-24), the asserted result is that the irrationality exponent of $\pi$ equals two, which also implies convergence of the classical Flint--Hills series.

\subsection{Mathematical statement of the claimed result}
\paragraph{Manuscript 1.}
For every real $\nu>2$, the claimed integer threshold $Q(\nu)$ satisfies
\[
 |\pi-p/q|\ge q^{-\nu}\qquad(p\in\mathbb Z,\ q\ge Q(\nu)).
\]
Thus the irrationality exponent of $\pi$ is two. The threshold is not claimed effective, and this does not assert a fixed positive lower bound $c/q^2$. The cited direct consequence is
\[
 \sum_{n=1}^{\infty}\frac1{n^3\sin^2 n}<\infty,
\]
with angles in radians.

\subsection{Outline of the argument}
For manuscript 1 (family 017; 2026-09-24), an interpolation determinant uses separated coordinate weights that remain valid independently of the chosen points. A full-rank minor has a lower arithmetic size, while repeated Taylor degrees give a quadratic analytic saving; a local multiplicity comparison turns the contradiction into the optimal approximation exponent.

\subsection{Scope relative to this problem}
The irrationality-exponent manuscript explicitly concludes convergence of the exact Flint Hills series displayed in this question, with angles in radians. Its ineffective threshold concerns the auxiliary approximation bound for $\pi$ and does not weaken the asserted convergence.

\subsection{Formal verification}
For manuscript 1 (family 017; 2026-09-24), the documented Lean scope is the approximation exponent of $\pi$, including its exact value two; the Flint--Hills-series consequence is absent from the selected statement.

This coverage does not establish a formal proof of the complete problem as stated here.

The linked family Comparator checks are \texttt{PiExponent.lean}.

Their named declarations include \texttt{OAI.PiExponent.main}.

This site has not executed the Lean code. Described statement coverage is not an independent proof check.

\subsection{Status caveat}
The manuscripts were produced by an internal OpenAI model. OpenAI cautions that results without formalization can contain mistakes. This site has not independently checked the proofs; the label records the scope of OpenAI's claim.

\subsection{Sources}
\begin{itemize}
\item OpenAI, \emph{The irrationality exponent of \(\pi\) is 2}, 2026-09-24; citation key \texttt{OAI:The-irrationality-exponent-of-pi-is-2-September-24-2026}. \href{https://github.com/openai/math/blob/adc7f1241b42e322a6451854ab7e4b4c146bf78a/preprints/The-irrationality-exponent-of-pi-is-2-September-24-2026/paper.pdf}{PDF (GitHub)}; \href{https://github.com/openai/math/raw/adc7f1241b42e322a6451854ab7e4b4c146bf78a/preprints/The-irrationality-exponent-of-pi-is-2-September-24-2026/paper.pdf}{PDF (direct)}; \href{https://github.com/openai/math/blob/adc7f1241b42e322a6451854ab7e4b4c146bf78a/preprints/The-irrationality-exponent-of-pi-is-2-September-24-2026/README.md}{Manuscript page and citation}; \href{https://github.com/openai/math/blob/main/preprints/The-irrationality-exponent-of-pi-is-2-September-24-2026/paper.pdf}{Latest version}; Statement source: \href{https://github.com/openai/math/blob/adc7f1241b42e322a6451854ab7e4b4c146bf78a/preprints/The-irrationality-exponent-of-pi-is-2-September-24-2026/build/main.tex\#L90}{thm:flinthills}; Statement source: \href{https://github.com/openai/math/blob/adc7f1241b42e322a6451854ab7e4b4c146bf78a/preprints/The-irrationality-exponent-of-pi-is-2-September-24-2026/build/main.tex\#L64}{thm:main}.
\item Family 017: \href{https://github.com/openai/math/blob/adc7f1241b42e322a6451854ab7e4b4c146bf78a/lean/docs/017.md}{Lean scope}.
\item \href{https://github.com/openai/math/blob/adc7f1241b42e322a6451854ab7e4b4c146bf78a/lean/ComparatorChallenges/PiExponent.json}{Comparator PiExponent.json}.
\item \href{https://github.com/openai/math/blob/adc7f1241b42e322a6451854ab7e4b4c146bf78a/lean/ComparatorChallenges/PiExponent.lean}{Comparator PiExponent.lean}.
\item Family 017: \href{https://github.com/openai/math/blob/adc7f1241b42e322a6451854ab7e4b4c146bf78a/reasoning\_traces/irrationality-exponent-of-pi.pdf}{Reasoning summary}.
\item \href{https://github.com/openai/math/blob/adc7f1241b42e322a6451854ab7e4b4c146bf78a/CONTENTS.md}{OpenAI catalogue entry}.
\end{itemize}
\end{document}
